\section{Introduction}

\IEEEPARstart{L}{ithium-ion} batteries are widely used in electric vehicles (EVs) and energy storage systems because of their high energy density, long cycle life, and low self-discharge rate \cite{1, 2}. In practical applications, a single cell cannot meet the voltage and capacity requirements of large-scale battery systems \cite{3, 4}. Battery packs therefore consist of many individual cells connected in series and parallel. Differences in manufacturing and operating conditions lead to variations in cell states and aging characteristics, which can reduce the usable capacity, charging efficiency, and durability of the battery pack and may also increase safety risks \cite{5, 6}. Effective charging control must therefore achieve fast charging while maintaining state-of-charge (SOC) balance among the cells and satisfying operating and safety constraints.

Most existing balancing strategies combine an equalization topology with a balancing algorithm. The equalization topology provides the physical path for transferring or dissipating energy, while the balancing algorithm determines how the topology is controlled to reduce differences among cells \cite{7, 8}. Their combined design affects balancing performance, energy efficiency, system complexity, and implementation cost \cite{ma2024two}.

\noindent\textit{Equalization topologies:}
Equalization topologies are generally divided into active and passive types. Active equalization transfers excess energy from cells with higher energy levels to cells with lower energy levels through switched capacitors, inductors \cite{9}, transformers \cite{10}, or DC-DC converters \cite{11}. This approach improves energy utilization, but the required circuitry becomes increasingly complex as the number of cells and series-parallel connections increases. Passive equalization instead uses discharge resistors connected in parallel with individual cells \cite{12}. Excess energy is dissipated as heat, which simplifies the circuit but reduces energy efficiency. Thus, conventional active and passive equalization both require additional hardware paths, resulting in either increased circuit complexity or direct energy loss.

\noindent\textit{Balancing algorithms:}
Traditional balancing algorithms are often based on constant-current constant-voltage (CC-CV) charging. Aizpuru {\it et al.} \cite{13} combined CC-CV charging with single-cell experimental data to construct a passive balancing system. Hoque {\it et al.} \cite{14} used particle swarm optimization (PSO) to optimize a CC-CV balancing controller. These methods are relatively simple to implement, but their ability to respond to changes in cell states and operating conditions is limited, and conservative constraints are often required to maintain safe operation.

More advanced control methods have also been studied for battery balancing. In \cite{15}, fuzzy logic control was used to adjust the equalization process adaptively. Fan {\it et al.} \cite{16} proposed a fast active balancing method based on model predictive control. These methods improve control flexibility and balancing performance, but their performance still depends on predefined control rules or models, which can limit their ability to accommodate changes in battery conditions during charging.

Deep reinforcement learning (DRL) offers another way to adapt charging decisions to the observed battery state. Yang {\it et al.} \cite{17} proposed a DRL-based framework for balancing-aware fast charging, while \cite{18} applied DRL to active equalization of redundant battery systems. By learning from interactions with the battery environment, these methods can adjust charging decisions as the battery state changes. However, two limitations remain. First, the reported DRL methods still rely on an equalization topology to redistribute energy among cells, and therefore retain the additional hardware and energy-transfer paths associated with conventional balancing systems. Second, reaching SOC balance at one point during charging does not guarantee that the cells will remain balanced afterward. Changes in cell states during continued charging may cause SOC differences to increase again.

Multi-agent reinforcement learning (MARL) provides a suitable framework for problems in which multiple decision makers act locally while sharing a common objective \cite{19, 20}. Liang {\it et al.} \cite{21} used MARL for real-time management of a battery swapping and charging system, while Dong {\it et al.} \cite{22} applied MARL to peak-shaving control. The battery charging problem considered in this work has a similar cooperative structure. Each cell has its own voltage, temperature, and SOC and can be assigned an independently regulated charging current, while all cells share the objectives of fast charging and SOC balancing. This makes it possible to coordinate the charging currents directly through MARL rather than relying on a separate equalization topology.

Based on this formulation, this paper proposes a multi-balancing-agent reinforcement learning (MBA-RL) framework for lithium-ion battery charging control. Each agent regulates the charging current of one cell, and the agents are coordinated through a shared battery-pack objective. A model-based safety mechanism is also introduced to evaluate the requested charging currents before execution. The framework is evaluated using cell-specific electrochemical-thermal models identified from cells with different cycling histories.

The main contributions of this paper are summarized as follows:

\begin{enumerate}

    \item An MBA-RL framework is developed to coordinate independently regulated cell charging currents without using a dedicated equalization topology. The cooperative control policy jointly considers charging progress and SOC balancing at the battery-pack level.

    \item A model-based safety mechanism is integrated into the MBA-RL framework. Before execution, the requested charging currents are checked against predicted voltage, temperature, and SOC constraints so that the charging policy operates within the prescribed cell-level limits.

    \item The proposed framework is evaluated using cell-specific electrochemical-thermal models identified from cells with different cycling histories. The evaluation covers unseen initial SOC conditions, operating variations, component ablations, and different battery-pack sizes.

\end{enumerate}

The remainder of this paper is organized as follows. Section II formulates the charging control problem. Section III presents the MBA-RL framework and the safety mechanism. Section IV reports and discusses the experimental results. Finally, Section V concludes this paper.